\section{Background}

\subsection{Kernel Attention}

Standard softmax attention:

\[
\text{Attn}(Q, K, V) = \frac{\exp(QK^\top / \sqrt{d_h}) V}{\exp(QK^\top / \sqrt{d_h}) \mathbf{1}}
\tag{3.1}
\]

If we can decompose the kernel $\kappa(q, k) = \exp(q^\top k / \sqrt{d_h})$ 
as $\phi(q)^\top \phi(k)$, then:

\[
\text{Attn}(Q, K, V)_i \approx \frac{\phi(q_i)^\top \sum_j \phi(k_j) v_j^\top}{\phi(q_i)^\top \sum_j \phi(k_j)}
\tag{3.2}
\]

The sums $\sum_j \phi(k_j) v_j^\top$ and $\sum_j \phi(k_j)$ are computed 
in $\mathcal{O}(L F d_h)$ and $\mathcal{O}(L F)$ respectively, giving 
overall \textbf{linear} complexity.

\subsection{Random Fourier Features}

\textbf{Bochner's theorem.} A continuous shift-invariant kernel 
$\kappa(x, y) = \kappa(x - y)$ is positive semidefinite iff it is the 
Fourier transform of a non-negative measure $p(\omega)$:

\[
\kappa(x - y) = \int p(\omega) e^{i \omega^\top (x - y)} d\omega
\tag{3.3}
\]

For the Gaussian kernel $\kappa(\Delta) = \exp(-\|\Delta\|^2 / 2)$:

\[
p(\omega) = (2\pi)^{-d/2} \exp(-\|\omega\|^2 / 2) = \mathcal{N}(0, I)
\tag{3.4}
\]

\textbf{Estimator.} Sampling $\omega_1, \dots, \omega_F \sim \mathcal{N}(0, I)$:

\[
\kappa(x - y) \approx \frac{1}{F} \sum_{i=1}^F \cos(\omega_i^\top (x - y))
= \frac{1}{F} \sum_{i=1}^F \left[\cos(\omega_i^\top x)\cos(\omega_i^\top y) + \sin(\omega_i^\top x)\sin(\omega_i^\top y)\right]
\tag{3.5}
\]

\textbf{Feature map.} Define:

\[
\phi(x) = \frac{1}{\sqrt{F}} [\cos(\omega_1^\top x), \dots, \cos(\omega_F^\top x), \sin(\omega_1^\top x), \dots, \sin(\omega_F^\top x)]
\tag{3.6}
\]

Then $\kappa(x - y) \approx \phi(x)^\top \phi(y)$.

\subsection{QR-Orthogonalization}

Sampling $\omega_i \sim \mathcal{N}(0, I)$ independently gives an unbiased 
estimator with variance $\mathcal{O}(1/F)$. A common trick is to 
\textbf{orthogonalize} the sampled vectors:

\[
[\omega_1, \dots, \omega_F] = Q \in \mathbb{R}^{d \times F}, \quad Q^\top Q = I_F
\tag{3.7}
\]

This reduces variance (Yu et al., 2016). We use QR:

\[
G \sim \mathcal{N}(0, I_{d \times F}), \quad G = Q R, \quad \omega_i = Q_{:, i}
\tag{3.8}
\]

\subsection{Causality in Linear Attention}

For autoregressive models, the sum $\sum_j \phi(k_j) v_j^\top$ must be 
causal: at position $i$, only $j \leq i$. This is achieved with a 
cumulative sum:

\[
\text{KV}_i = \sum_{j \leq i} \phi(k_j) v_j^\top, \quad
\text{Z}_i = \sum_{j \leq i} \phi(k_j)
\tag{3.9}
\]

Both are computed in $\mathcal{O}(L)$ via \texttt{cumsum}.

\subsection{LayerNorm}

LayerNorm normalizes per example:

\[
\text{LN}(x) = \gamma \odot \frac{x - \mu}{\sigma + \epsilon} + \beta
\tag{3.10}
\]

We apply LN on the \textbf{head dimension} before the random projection.

\subsection{Gated Residual}

A scalar gate $\alpha$ blends two branches:

\[
\text{out} = \alpha A + (1 - \alpha) x
\tag{3.11}
\]

Initialized to $\alpha_0 = 0.88$ (via $w_0 = 2$), which favors attention.

\subsection{Notation}

\begin{table}[H]
\centering
\small
\begin{tabular}{ll}
\toprule
\textbf{Symbol} & \textbf{Meaning} \\
\midrule
$\kappa$ & Kernel function \\
$p(\omega)$ & Spectral distribution \\
$F$ & Feature dimension \\
$\phi$ & Feature map \\
$\omega_i$ & Random frequency \\
$W$ & Projection matrix \\
$\alpha$ & Preservation gate \\
\bottomrule
\end{tabular}
\end{table}

